Nonlinear physical dynamics, such as those governing cislunar aerospace applications, frequently distort initial uncertainties into complex, non-Gaussian probability distributions. Direct evaluation and filtering within these intractable target state spaces are computationally prohibitive. Optimal transport provides a theoretical mechanism to construct a deterministic coupling that pushes these complex target distributions forward to a simple, standard normal reference distribution, rendering subsequent analytical tasks mathematically tractable. However, guaranteeing that these transport maps remain physically valid and invertible requires strict adherence to monotonicity constraints, which inherently restricts the operational domain of the learning algorithm to a bounded feasible region.

In this chapter, we formulate a deterministic density estimation framework utilising the lower-triangular Knothe-Rosenblatt rearrangement. By parameterising the map components with a truncated probabilist's Hermite polynomial basis, we reformulate the infinite-dimensional map estimation task into a finite-dimensional, constrained maximum likelihood problem. The physical necessity of probability mass conservation dictates a dense system of linear inequality constraints evaluated over an empirical particle ensemble. This structural formulation forms the applied aerospace pillar of this thesis, directly establishing the exact polyhedral geometry and Constrained Quadratic Program (CQP) that necessitates the accelerated Dykstra projection algorithm developed in the preceding chapter.

\section*{Notation}
We denote the continuous state space as $\mathcal{X} \subseteq \mathbb{R}^d$, where $d$ is the spatial dimension. A discrete ensemble of particles drawn from the target measure is denoted by $x^{(i)}$, with the total empirical particle count fixed as $n$. The deterministic transport map is denoted as $S(x)$, and its individual scalar components are indicated by a superscript as $S^k(x)$. The partial derivative with respect to the $k$-th spatial coordinate is denoted by $\partial_k$. The probability density functions of the target and reference measures are denoted as $p(x)$ and $\eta(z)$, respectively. The univariate probabilist's Hermite polynomial of degree $m$ is denoted as $\mathrm{He}_m(\cdot)$, while the multivariate basis functions constructed via tensor products are denoted as $\Psi_j(x)$. The unknown scalar coefficients comprising the decision variable for the $k$-th map component are grouped in the vector $w^{(k)}$. The dense constraint matrix is denoted as $A$, and the resulting feasible polyhedral set enforcing physical monotonicity is denoted as $\mathcal{C}$.